% 原书第 8.5 节。
\section{大语言模型选型}
\label{sec:ch08-selection}
为应用选出最合适的大语言模型，可以说是其中最重要的一项决策。语言模型及其建模方法的能力不断提高，OpenAI 在 2022 年 11 月发布 ChatGPT 是一个显著节点。这是撰写本书的主要动因，也使采用这些模型的领域涌现出大量创新。不过，选用哪一种模型时，模型能力或分析质量只是一项考量。

\begin{keypoint}
\textbf{为模型选型和定制方式绘制决策流程图}

以下各节讨论大语言模型的各种属性及其与应用开发情境的关系。读者可以自行绘制决策流程图，把模型属性对应到应用领域。这样做有助于解决应用所针对的问题，使方案日后仍可维护，并在提高效率和效果的同时，减少采用大语言模型带来的风险。请仔细考虑应用的功能需求和非功能需求，将其与下文讨论的机会和难题逐项联系起来，据此建立决策过程。
\end{keypoint}

模型选择发生在项目早期，会影响许多后续选项，因此还应考虑许多其他标准和属性。图~\ref{fig:ch08-selection}改绘自 \refcite{ch08-ref-047}{Li et al. (2023)}，以一个领域的选型和开发过程为例，说明如何在开源和闭源模型之间作出选择，并根据工具、数据和预算等条件选择定制路径。定制方法自上而下按成本由低到高排列，便于在考虑开销时依次尝试。

\begin{figure}[!htbp]
\centering
\begin{tikzpicture}[
  x=1cm,y=1cm,font=\sffamily\footnotesize,
  condition/.style={ellipse,draw=cyan!50!black,fill=cyan!13,align=center,
    text width=2.33cm,minimum width=3.3cm,minimum height=.78cm,inner sep=1pt},
  choice/.style={rectangle,rounded corners=2pt,draw=orange!65!black,fill=orange!15,
    align=center,text width=2.57cm,minimum height=.83cm,inner sep=3pt},
  level/.style={align=center,text width=1.16cm,font=\sffamily\scriptsize,inner sep=0pt},
  yes/.style={-{Stealth[length=1.7mm]},draw=green!50!black},
  no/.style={-{Stealth[length=1.7mm]},draw=red!70!black},
  next/.style={-{Stealth[length=1.7mm]},draw=black!80},
  route/.style={font=\sffamily\scriptsize,fill=white,inner sep=1pt}]
\node[condition] (private) at (0,0) {数据需要\\保密吗？};
\node[condition] (examples) at (0,-1.78) {能提供若干\\任务示例吗？};
\node[condition] (tools) at (0,-3.56) {有工具可辅助\\注入知识吗？};
\node[condition] (domain) at (0,-5.34) {有领域数据\\可用于微调吗？};
\node[condition] (budget) at (0,-7.28) {预算逾 100 万\\美元且有数十亿\\训练词元吗？};
\node[choice] (own) at (-3.40,-.79) {1A 自行部署\\开源大语言模型};
\node[choice] (api) at (3.40,-.79) {1B 使用第三方\\大语言模型（如 GPT-4）};
\node[choice] (few) at (-3.40,-2.57) {2A 对选定模型\\采用少样本方法};
\node[choice] (rag) at (-3.40,-4.35) {3A 用工具增强\\所选大语言模型};
\node[choice] (tune) at (-3.40,-6.13) {3B 通过微调\\向模型注入知识};
\node[choice] (train) at (-3.40,-8.25) {4A 从头训练\\自己的大语言模型};
\node[draw,rounded corners,fill=black,text=white,font=\sffamily\scriptsize] (end) at (3.40,-8.25) {无可行方案};
\draw[yes] (private.west) to[out=180,in=90] node[route,above,pos=.52]{是} (own.north);
\draw[no] (private.east) to[out=0,in=90] node[route,above,pos=.52]{否} (api.north);
\draw[no] (private) -- node[route,right]{否} (examples);
\draw[next] (own.south) to[out=-90,in=180] (examples.north west);
\draw[next] (api.south) to[out=-90,in=0] (examples.east);
\foreach \question/\option/\following in {examples/few/tools,tools/rag/domain,domain/tune/budget}{
  \draw[yes] (\question.south west) to[out=180,in=90] node[route,above,pos=.52]{是} (\option.north);
  \draw[no] (\question) -- node[route,right]{否} (\following);
  \draw[next] (\option.south) to[out=-90,in=180] (\following.north west);
}
\draw[yes] (budget.south west) to[out=180,in=90] node[route,above,pos=.52]{是} (train.north);
\draw[no] (budget.east) to[out=0,in=90] node[route,above,pos=.52]{否} (end.north);
\draw[next] (train.east) -- (end.west);
\node[level] at (-5.42,-.79) {第 1 级\\零样本};
\node[level] at (-5.42,-2.57) {第 2 级\\少样本};
\node[level] at (-5.42,-5.24) {第 3 级\\RAG 或\\微调};
\node[level] at (-5.42,-8.25) {第 4 级\\从头训练};
\draw[next] (2.12,-2.82) -- (2.82,-2.82) node[right,font=\sffamily\scriptsize]{效果不满意};
\node[font=\sffamily\scriptsize,draw=orange!65!black,fill=orange!15,rounded corners=2pt]
  at (3.29,-3.47) {决策框};
\node[font=\sffamily\scriptsize,ellipse,draw=cyan!50!black,fill=cyan!13]
  at (3.29,-4.17) {方案框};
\draw[-{Stealth[length=3mm]},line width=1.25pt,draw=orange!75!black]
  (3.29,-4.77) -- node[right,font=\sffamily\scriptsize]{成本上升} (3.29,-6.85);
\end{tikzpicture}
\caption{金融服务领域中开源与闭源大语言模型的选型及定制流程。}
\label{fig:ch08-selection}
\end{figure}


\begin{keypoint}
为项目选型时，总拥有成本（total cost of ownership，TCO）也是一种有用的分析框架。它按项目的模型、使用场景等细节估算各项成本，再加总以便比较方案。部分项目如下。
\begin{itemize}
\item \textbf{每词元成本}，反映每次查询的处理和生成开销。
\item \textbf{人力成本}，估算构建和部署大语言模型服务所需的人力开销。
\item \textbf{整体搭建成本}，估算部署、运行和管理大语言模型服务的总开销。
\end{itemize}
希望进一步了解这些因素及计算过程的开发人员，可以从 Hugging Face 的资料入手，其中包括一个交互式 TCO 计算器。读者也可以查阅类似资料，了解这一框架怎样辅助应用开发中的决策。
\end{keypoint}


\transnote{所据英文版标出了注 3 和注 4 的序号，页面上未印出这两条注释的链接。}

各类大语言模型各有利弊，哪些利弊重要很大程度上取决于应用情境。例如，一个应用主要做文本摘要，选用号称在实体抽取基准上达到当时最佳水平的模型，意义不大。本节介绍选型的主要标准，帮助开发人员判断模型属性是否切合其应用，从而作出选择。

\subsection{开源与闭源}
开源还是闭源，是开发人员为应用选择大语言模型时首先考虑的问题之一。两者的主要取舍通常在使用灵活性、便利性和成本三个方面，还应考虑表~\ref{tab:ch08-open-closed}所列的其他因素。例如，开发团队起初可能看中开源模型的使用成本低、灵活性高，但按表中各项标准逐一评估后，有可能发现采用开源模型的应用，其总拥有成本远高于采用闭源模型的应用。因此，开发人员应尽量全面地评估模型选择。以下各项分别讨论开源和闭源模型的有关倾向。

{\small\setlength{\tabcolsep}{4pt}\renewcommand{\arraystretch}{1.12}
\begin{longtable}{@{}>{\raggedright\arraybackslash}p{15mm} >{\raggedright\arraybackslash}p{47mm} >{\raggedright\arraybackslash}p{47mm}@{}}
\caption{应用开发中选择开源或闭源大语言模型时应考虑的方面。}\label{tab:ch08-open-closed}\\
\toprule
方面 & 闭源 & 开源\\
\midrule
\endfirsthead
\multicolumn{3}{@{}l}{表 \thetable\ （续）}\\
\toprule
方面 & 闭源 & 开源\\
\midrule
\endhead
\midrule
\multicolumn{3}{r@{}}{下页续}\\
\endfoot
\bottomrule
\endlastfoot
可获取性 & 终端用户往往无法获取模型架构、训练数据、学习目标及源代码的详细信息。 & 终端用户可依宽松许可获取模型架构、训练数据、学习目标和源代码的详细信息。\\
透明度 & 终端用户通常通过预先规定的 API 使用模型，上述各项的透明度有限，影响在不同情境中的使用，也影响应用的调试、维护和支持。 & 终端用户通常可获得模型权重。开发人员负责托管模型并提供服务。透明度高，支持多种定制和灵活的使用场景，相应支持也由开发人员承担。\\
成本 & 使用成本通常高于开源模型，但便利性和使用支持等因素往往可以抵消较高的价格。 & 使用成本可能远低于闭源模型，部分原因是终端用户可以定制并优化计算资源消耗等。较低的使用成本应与开发和支持开销一并权衡，后两项均由终端用户承担。\\
定制 & 仅限于规定的方法。这些方法的成本往往高于开源方案。 & 可灵活采用适应性预训练、全参数微调、参数高效微调或基于提示的学习。\\
使用方式 & 通常限于若干预先设定的方式，例如图形用户界面或 API。 & 使用方式灵活，可由终端用户自行设计和开发。\\
协作 & 仅限于已集成的协作工具。 & 协作方式灵活，可采用终端用户所选的任何协作工具或框架。\\
安全 & 由模型所有者全面管理。这对没有专门安全资源的小型开发团队很有价值。然而，任何安全方案都无法做到百分之百有效，闭源条件也可能使模型或服务中的漏洞难以发现。 & 终端用户必须自行开发和管理安全措施。这样可获得最大的灵活性，但对小型开发团队而言，承担这项责任可能复杂且昂贵。\\
隐私 & 输入模型的任何数据均受终端用户协议约束。提供方可能获准用这些数据改进未来版本，因此还要考虑知识产权流失。 & 终端用户根据需求或偏好决定隐私保护的程度。\\
支持 & 终端用户可获得模型所有者的支持，并负责为自己应用中封装的功能提供支持。 & 模型往往不附带终端用户支持。活跃的开源社区有时可通过 GitHub 或 Hugging Face 项目为用户解决问题提供帮助。\\
更新与维护 & 模型所有者安排并发布更新。更新过程可能不透明，并受模型所有者的商业规划支配，未必符合所有终端用户的意愿。 & 终端用户负责全部模型更新与维护。这样可获得最大的控制权，但对较小的开发团队而言，可能需要承担较高的开销。\\
\end{longtable}
}

\subsection{分析质量}
在模型选型的各项考量中，权重最高的通常是模型完成所要解决的使用场景相关任务的质量。较大的模型往往具有较强的分析表现，因而计算成本和财务开销是主要的取舍。汇总的排行榜可作为衡量分析质量的参考。排行榜用一系列标准化基准评估大语言模型，便于直接比较和筛选。解读结果时应仔细考虑使用的评估方法（\refcite{ch08-ref-003}{Alzahrani et al., 2024}）。

即便如此，排行榜也可以大致显示模型之间的相对表现，有助于迅速缩小候选范围，再进一步评估各模型是否适合项目。原书认为第~\ref{sec:ch08-categories}~节所述的 Hugging Face Open LLM Leaderboard 在当时维护得最好，但只覆盖开源模型。也有同时覆盖开源和闭源模型的排行榜，但这些项目能否长期维护并不明确，例如 \url{https://llm-leaderboard.streamlit.app/}。希望同时评估两类模型时，可能需要检索合适的资料。

把候选模型缩小到可管理的范围后，可以用更有针对性的任务进一步评估分析表现。大语言模型的任务表现对所用数据敏感。换用更接近应用领域的数据分布，模型的实际表现可能低于排行榜上的基准结果。第二轮贴近使用场景的评估可进一步缩小范围，找出在该评估中表现最好的模型，或在项目容许继续开展提示工程、预训练或微调时有改进潜力的模型（\refcite{ch08-ref-108}{Yuan et al., 2023}）。

\subsection{推理延迟}
\label{sec:ch08-latency}
大语言模型的推理延迟是选型的重要因素，可以理解为从处理单个请求到把对应响应返回给用户所需的总时间，见图~\ref{fig:ch08-latency}。把用户的输入提示传到模型 API 的延迟（图中第~1~步），以及把模型响应从 API 送回用户的延迟（第~4~步），大部分可通过网络优化来处理。除此之外，对总体延迟影响最大的有两个推理阶段，即模型网络处理输入提示中词元的时间（第~2~步），以及抽样和组织响应词元的时间（第~3~步）。前者称为预填充（prefill），后者称为解码（decode）（\refcite{ch08-ref-002}{Agrawal et al., 2024}）。

\sourcefigure{fig08_03.pdf}{大语言模型的延迟由四部分组成。（1）用户提示到达模型推理 API 的时间。（2）模型处理用户提示的时间。（3）模型抽取相关词元并组成响应的时间。（4）模型响应从 API 传给用户的时间。}{fig:ch08-latency}

\urldef\latencytableurl\url{https://huyenchip.com/2023/04/11/llm-engineering.html#cost_and_latency}
由于采用 Transformer 架构，预填充阶段可并行处理提示词元，因此延迟相对较高（与解码步骤相比），计算资源利用率也高。相较而言，解码是顺序过程，因为在生成下一个输出词元之前，必须先生成所有之前的词元。因此，单个输出词元的延迟相对较低，顺序生成也使计算资源利用率较低。这意味着与输出长度相比，提示中的输入词元数对推理延迟的影响不大。例如，表~\ref{tab:ch08-latency}\footnote{表中数据转载自 \latencytableurl。}给出了调整输入和输出词元长度对 OpenAI \texttt{gpt-3.5-turbo} 模型响应延迟的影响。把输入词元从 51 个增至 232 个，而输出词元维持为 1 个，延迟变化很小。输入长度大致相同时，若把输出从 1 个词元增至 26 个，延迟接近原来的 3 倍。这显示输入和输出长度对推理延迟的影响并不均衡。



\begin{table}[htbp]
\centering
\caption[输入和输出长度对推理延迟的影响]{输入和输出长度对推理延迟的影响。数据来自 OpenAI 的 \texttt{gpt-3.5-turbo} 模型。OpenAI 如何调度和路由用户查询，用户无从得知，因此部分波动来自 API 延迟。即便绝对延迟发生变化，输入和输出长度设置之间的关系仍保持一致。p50 延迟以秒计，表示在 $n=20$ 次请求中，有 50\% 的请求响应时间不超过该值。}
\label{tab:ch08-latency}
\begin{tabular}{rrr}
\toprule
输入词元数 & 输出词元数 & p50 延迟（秒）\\
\midrule
51 & 1 & 0.58\\
232 & 1 & 0.53\\
228 & 26 & 1.43\\
\bottomrule
\end{tabular}
\end{table}

考虑到这种不均衡，大语言模型有哪些属性会影响推理延迟？最直观的是模型规模。经验上，参数越多，延迟越高。参数更多的模型需要更多计算来处理输入并生成输出。模型架构也很重要。层数、各层的复杂程度、Transformer 块所用的注意力机制，以及 Transformer 块在网络中的数量和位置，都会影响推理延迟。

参数以何种数值精度存储，也是影响大语言模型推理延迟的重要因素。第~\ref{ch:adaptation}~章第~\ref{sec:ch04-quantization}~节已详细讨论量化。在比较开源和闭源模型时，两者在定制方面的差别值得注意。闭源模型定制的限制较多，终端用户只能采用模型所有者支持的量化方式。开源模型的终端用户通常可以自行测试并采用最适合使用场景的量化方法。量化有可能显著降低推理延迟，以及运行或托管模型所需的内存和存储开销，因此应认真考虑闭源模型在定制上的限制。如果预计请求与响应的次数很少，这一点可能不太重要。然而，在请求与响应次数较多的情况下，闭源大语言模型可能成为应用的瓶颈。例如，按原书当时的描述，OpenAI 的 API 通常按端点和模型分别设定速率限制。

\subsection{成本}
应用开发使用大语言模型时，很多环节都要付出成本。人们往往只考虑推理的每词元成本，它确实是最重要的成本之一。不过，大量研究和商业投入不断压低这项成本，其数值也随之改变。例如，图~\ref{fig:ch08-costs}给出了当时的情况。\footnote{转载自 \url{https://medium.com/@ja_adimi/comparison-cost-analysis-should-we-invest-in-open-source-or-closed-source-llms-bfd646ae1f74}。}围绕开源与闭源模型的讨论通常认为，闭源模型每个词元的成本显著高于开源模型。不过，这很可能只适用于能力最强的闭源模型版本。原书以 OpenAI 当时的定价为例，较旧型号的推理开销通常只是最新型号的一小部分。

\sourcefigure{fig08_04.pdf}{生成 100 万个词元的美元成本。图中比较 OpenAI 的两个闭源大语言模型 GPT-4 和 GPT-3.5 Turbo，以及在 A100 和 A10G 两种 GPU 上运行的开源模型 Llama-2-7B。两种配置的 Llama-2 生成成本均低于 OpenAI 当时能力最强的模型 GPT-4 的十分之一，但 GPT-3.5 Turbo 比这两种配置都便宜。选择项目所用模型时，应考虑闭源模型推理成本下降的趋势。费用为 2023 年 8 月分析时的数值。}{fig:ch08-costs}

采用总拥有成本框架时，除了每词元生成成本，开发人员也应估算应用开发生命周期各环节的开销。例如，采用开源模型的推理成本可能更低，分析质量也可能较低。微调或投入提示工程也许能缩小这种差距，但优化本身要付出成本。微调或测试都需要数据，收集、标注和整理数据可能费力且复杂。此外，模型经微调后还需要长期维持其表现，这意味着持续评估并在必要时改进模型的工作也应计入成本。定制和维护大语言模型是复杂的技术工作，因而相较于采用闭源模型，项目需要承担更高的专业技术与人才成本。选择开源模型时，托管和计算资源管理也是直接开销，会增加整体复杂度和成本。

\subsection{适应性与维护}
开源大语言模型通常更易适配，因为终端用户往往可以直接获得模型权重、训练数据和源代码。这使用户可以采用第~3~章和第~4~章介绍的一项或多项技术来适配或定制模型，进而控制应用的行为和表现。灵活性的代价是开发生命周期需要消耗更多资源。

闭源模型适应性较低，相应地，其开发周期消耗的资源也少得多。如果项目利用大语言模型执行常见任务，专有模型很可能已经具备较好的开箱能力，无需再适配或定制。类似地，采用 $n$ 样本上下文学习等高级提示工程技术，还可以进一步改善结果。这些情况下可以考虑闭源或专有模型。持续维护仍然是决策因素。闭源模型的维护通常不透明，会按照模型所有者的规划或维护日程进行。假如变更可以不事先通知终端用户，开发人员就应了解所选模型悄然升级对应用表现的影响。升级后的模型会降低用户体验吗？如果处置不当，是否会引发伦理或安全风险？

通过在开发生命周期中引入持续监控和评估，可以减少相当一部分这类风险。不过，大语言模型的规模和复杂性使得事先预测所有失效情形成为不可能。因此，选用何种大语言模型，取决于具体使用场景。在应用出现错误的代价很高的环境中（例如受监管行业），可能无法任由用户的结果碰运气，或由 OpenAI、Anthropic 等第三方决定。此时，应用提供方只能选择自己能够全面负责这些风险的模型。

\subsection{数据安全与许可}
应用常借助来自用户或其他来源的敏感数据提供结果。如果把敏感数据作为上下文写进提示并发送给大语言模型，很难彻底阻止模型响应包含敏感上下文，因此会出现数据安全或隐私风险。许多方案采用先生成再过滤的方法，在向用户提供响应之前删除其中的敏感数据。也可采用第~5~章综述的对齐方法，减少用户收到敏感或不希望出现的信息的风险。

对于闭源模型或由第三方托管的模型，数据经网络传送到模型 API 会带来额外的安全问题。如果模型为专有模型，还须考虑所有者如何留存和使用数据。ChatGPT 刚向公众开放时，OpenAI 如何利用服务中突然涌入的用户查询和交互，是一个热门话题。许多用户担心 OpenAI 会把他们的提示用于进一步改进 ChatGPT，其中可能包含私密或敏感信息，进而可能暴露给其他用户。由于 ChatGPT 的流行程度增长很快，这个问题当时并不明朗。关于通过越狱从大语言模型提取“训练数据”的各种报告，使医疗和金融服务等对数据隐私及安全要求很高的行业中，许多潜在用户心存疑虑（\refcite{ch08-ref-107}{Yu et al., 2024}）。

规划应用设计及维护工作时，开发人员应了解许可方面的两个问题。其一与传统软件和工具的许可相似，即大语言模型发布时所用的许可往往规定了应用范围。例如，某些许可只准用于研究或实验，另一些可能使应用的营销或推广比预期复杂。本节不详细讨论不同的许可类型及其对应用开发人员的影响，感兴趣的读者可查阅有关资料。\footnote{\url{https://github.com/eugeneyan/open-llms}}

其二是预训练数据的许可。越来越多的材料表明，至少部分模型训练数据可能受版权限制，把这些材料纳入训练集可能构成侵权。备受关注的案例之一，是美国法院审理的 \textit{The New York Times} 与 OpenAI 之间的案件。在这起案件中，《纽约时报》主张 OpenAI 未经授权将该媒体的内容纳入大语言模型，侵犯了其权利（\refcite{ch08-ref-087}{S.D.N.Y., 2023}）。在此类案件的审理过程中，应用开发人员需预判裁决对应用维护的影响。如果法院支持著作权主张方，模型所有者可能须撤下使用受版权保护的材料训练的模型，并发布更新版本。应用开发人员因而可能不得不更换应用中的模型，并补做测试以确认表现是否一致。
